% !TeX root = ../main.tex
\chapter{System Implementation}
\label{chap:system_implementation}

\section{Introduction}
\label{sec:system_implementation_intro}
As discussed in Chapter~\ref{chap:system_architecture}, the proposed solution follows an end-to-end streaming architecture composed of multiple interconnected layers. 

While the previous chapter focused on the architectural design decisions and the rationale behind the selected technologies, this chapter presents the practical implementation of the system. First, Section~\ref{sec:implementation_methodology} introduces the implementation methodology, version control practices, and continuous integration workflows. Subsequently, the implementation details of each architectural layer are described, together with the deployment strategy, service integration mechanisms, and the technologies used to establish communication between components. 

Finally, the chapter explains how the different layers are combined to form a complete data pipeline, from telemetry generation at the source systems to data storage and downstream analytics.

\section{Implementation Methodology}
\label{sec:implementation_methodology}
\textit{Architectural counterpart: Chapter~\ref{chap:system_architecture}, Section~\ref{sec:arch_decisions}.}

Translating the high-level architectural design presented in Chapter~\ref{chap:system_architecture} into an operational platform requires establishing a structured implementation methodology. The proposed solution constitutes a distributed and heterogeneous system that spans physical vehicle simulation engines written in Rust, IoT edge messaging brokers, distributed stream processors powered by Apache Spark, cloud-native lakehouse storage tiers, and deep learning models built in PyTorch. Managing such a multi-technology platform demands disciplined software engineering and DevOps practices to guarantee reproducibility, traceability, and seamless interoperability between components.

\subsection{Distributed Version Control and Commit Traceability}
Modern software engineering relies on distributed version control to record changes, maintain history, and coordinate technical iterations. In this project, Git~\cite{git} was selected as the central version control system, hosted within the Bitbucket~\cite{bitbucket} repository infrastructure of the industrial partner.

\subsection{Continuous Integration and Automated Verification}
Continuous Integration (CI) and Continuous Deployment (CD) principles were integrated into the development workflow to ensure that modifications to individual microservices do not destabilize the overall pipeline. The containerized design of the platform directly supports automated verification workflows:
\begin{itemize}
    \item \textbf{Automated Build Validation:} Because system components are containerized, automated build steps verify that Docker images can be built without errors, ensuring that third-party dependencies and compilation targets (such as the Rust-based vehicle simulator) remain consistent and pinned.
    \item \textbf{Quality Assurance and Static Analysis:} Static code analysis, syntax checking, and formatting tools were applied across the various languages utilized in the project (Python, Rust, shell scripts, and configuration files) to enforce consistent code quality standards.
    \item MORE WILL BE ADDED!
\end{itemize}

\section{Containerized Deployment}
\label{sec:containerized_deployment}
\textit{Architectural counterpart: Chapter~\ref{chap:system_architecture}, Section~\ref{sec:arch_decisions}.}
Following the architectural decisions described in Chapter~\ref{chap:system_architecture}, all components were deployed as Docker containers. Containerization provides reproducibility, isolation of dependencies, and simplifies the deployment and maintenance of the platform. 

\subsection{Docker Compose Organization} 
To improve maintainability and modularity, the platform is organized as a collection of independent Docker Compose files. Each architectural layer maintains its own deployment definition, while a main \texttt{docker-compose.yml} file aggregates all services using the Docker Compose \texttt{include} directive. 

This approach allows each subsystem to be developed and deployed independently while preserving a unified deployment strategy for the complete platform. 

Listing~\ref{lst:main-compose} shows the complete root Docker Compose configuration used to orchestrate all subsystems, define CA certificates, and declare shared network boundaries.

\begin{lstlisting}[language=yaml, caption={Root Docker Compose configuration}, label={lst:main-compose}, basicstyle=\ttfamily\small, breaklines=true]
secrets:
  ca_cert:
    file: ./certs/edag-r1.crt

include:
  - docker/mqtt-broker/docker-compose.yml
  - docker/source_systems/vda5050_vehicle_simulator/docker-compose.yml
  - docker/redpanda/docker-compose.yml
  - docker/source_systems/telegraf/docker-compose.yml
  - docker/data-storage/questdb/docker-compose.yml
  - docker/cross-cutting/docker-compose.yml
  - docker/utils/sim-control-panel/docker-compose.yml
  - docker/stream-processing/spark/docker-compose.yml
  - docker/data-lakehouse/minio/docker-compose.yml

networks:
  lakehouse-net:
    driver: bridge
  redpanda-net:
    driver: bridge
  mqtt-net:
    driver: bridge
  questdb-net:
    driver: bridge
\end{lstlisting}

\subsection{Network Segmentation} 
Dedicated Docker bridge networks were created to isolate communication between system components. Rather than exposing every service to a common network, individual networks were assigned according to the communication requirements of each layer. 

This strategy reduces unnecessary connectivity between containers and simplifies service discovery. Furthermore, it follows the principle of least privilege by restricting access only to the services required by each component. 

The platform defines four segmented bridge networks:
\begin{itemize}
    \item \texttt{mqtt-net}: Dedicated to edge communication between simulated AGVs, the Eclipse Mosquitto broker, and the Telegraf ingestion agent.
    \item \texttt{redpanda-net}: Connects the central distributed event streaming brokers, Telegraf producer, Spark streaming consumers, and Kafka Connect.
    \item \texttt{lakehouse-net}: Dedicated to the analytical storage tier, connecting Spark stream processing with the MinIO S3 object store.
    \item \texttt{questdb-net}: Isolates the operational time-series database, schema initialization containers, Kafka Connect sink, and Grafana monitoring dashboards.
\end{itemize}

\subsection{Modular Service Configuration}\label{subsec:modular-service-configuration}
To improve modularity, maintainability and deployment portability, all infrastructure dependencies are configured externally through environment variables.
Rather than hardcoding service endpoints, credentials or storage locations within the application source code, each component retrieves its configuration at runtime from the shared \texttt{.env} file. 
This design allows individual services to be replaced without requiring modifications to the application logic.
As a consequence, Spark jobs, Kafka clients, Iceberg catalogs and storage backends remain loosely coupled and can be adapted to different deployment environments through configuration only.  
This enables the same container images to be reused across development, testing and production deployments. 
For example, the current implementation uses Redpanda as the streaming platform: 

\begin{lstlisting}[language=bash] 
  KAFKA_BROKER=redpanda:9092 
\end{lstlisting} 

If the system were migrated to a different Kafka-compatible broker, only the endpoint definition would need to be updated: 

\begin{lstlisting}[language=bash] 
  KAFKA_BROKER=my-broker.example.com:9092 
\end{lstlisting} 

No changes to the Spark jobs, data processing logic or container images would be required.

\section{Source Systems: Vehicle Simulation and Physical Modeling} 
\label{sec:source_system_simulation}
\label{sec:physical_modeling} 
\textit{Architectural counterpart: Chapter~\ref{chap:system_architecture}, Section~\ref{sec:arch_source_systems}.}
This section describes the implementation of the source systems layer, the origin of truth of the platform. It covers the containerized deployment of the extended vehicle simulator, the electromechanical, kinematic, thermal and battery model that governs the generated telemetry, the physical fault-injection and transducer models, and the resulting telemetry channels published as VDA~5050 payloads.

\subsection{Simulator Architecture and Containerized Deployment}
\label{subsec:simulator_deployment}

The source system layer is currently implemented using a modified version of the \texttt{vda5050\_vehicle\_simulator} \cite{bostanciVda5050}. The simulator executes inside a dedicated container and connects to the MQTT broker through the shared \texttt{mqtt-net} network. 

The \texttt{depends\_on} directive combined with Docker health checks guarantees that the simulator starts only after the broker has become available. 

\begin{lstlisting}[language=yaml, caption={VDA5050 Simulator Docker Compose}, label={lst:vda5050-simulator-docker-compose}, basicstyle=\ttfamily\small, breaklines=true]
services:
  simulator:
    image: vda5050-vehicle-simulator:latest
    volumes:
      - ./config.toml:/app/config.toml:ro
    networks:
      - mqtt-net
    depends_on:
      mosquitto:
        condition: service_healthy
    restart: unless-stopped
\end{lstlisting}

The simulator configuration is loaded from a dedicated configuration file that defines fleet parameters, MQTT communication settings, and telemetry generation properties. During execution, each simulated AGV periodically publishes VDA 5050-compliant state messages to the MQTT broker. 

To support the anomaly detection use case, the original simulator was extended with additional telemetry channels representing motor, battery, and drivetrain behavior. These extensions enable the generation of realistic training data for machine learning workloads while preserving compatibility with the VDA 5050 communication standard.

To provide physically consistent telemetry for anomaly recognition, the kinematic simulation of the automated guided vehicle (AGV) is extended with an electromechanical powertrain model. The objective is not to replicate every microscopic effect of a physical robot, but to capture the dominant coupling between trajectory demands, mechanical loads, electrical responses, and thermal dissipation.

\subsection{Motor and Drivetrain Dynamics}

The traction drive is modeled as an equivalent permanent-magnet DC (PMDC) machine coupled to a reduction gearbox ($i_g$) and driven wheel ($r_w$).

The electrical armature circuit is governed by:
\begin{equation}
    V_m(t) = R_a(T_m) I_m(t) + L_a \frac{dI_m(t)}{dt} + K_e \omega_m(t)
    \label{eq:armature_base}
\end{equation}
where $V_m$ is the armature terminal voltage, $I_m$ the armature current, $R_a(T_m)$ the winding resistance, $K_e$ the back-EMF constant, and $\omega_m$ the rotor speed.

\paragraph{Quasi-static approximation ($L_a \approx 0$):}
Because the electrical time constant ($\tau_e = L_a / R_a \approx 40\,\mathrm{ms}$) is much shorter than the mechanical time constant ($\tau_m \approx 1\text{--}3\,\mathrm{s}$) and the state reporting rate, the inductive transient $L_a \frac{dI_m}{dt}$ is neglected. This yields the quasi-static relationship:
\begin{equation}
    V_m(t) = R_a(T_m) I_m(t) + K_e \omega_m(t)
\end{equation}

Electromagnetic torque follows:
\begin{equation}
    \tau_m(t) = K_t I_m(t)
    \label{eq:electromagnetic_torque}
\end{equation}
where $K_t$ is the torque constant. The mechanical torque balance at the motor shaft is:
\begin{equation}
    J_m \frac{d\omega_m(t)}{dt} + B_m \omega_m(t) = \tau_m(t) - \frac{\tau_L(t)}{\eta_g i_g} - \tau_{\mathrm{friction}}
    \label{eq:torque_balance}
\end{equation}
where $J_m$ is rotor inertia, $B_m$ viscous friction, $\eta_g$ gearbox efficiency, and $\tau_{\mathrm{friction}}$ unmodeled or degraded drivetrain friction. The wheel load torque $\tau_L(t)$ aggregates rolling resistance, road grade, vehicle acceleration, and drawbar forces:
\begin{equation}
    \tau_L(t) = r_w \left[ C_{rr} m_{\mathrm{total}} g \cos\theta + m_{\mathrm{total}} a(t) + m_{\mathrm{total}} g \sin\theta + F_{\mathrm{drawbar}} \right]
\end{equation}
where $m_{\mathrm{total}} = m_{\mathrm{chassis}} + m_{\mathrm{payload}}$ is the total vehicle mass.

\subsection{Kinematics and Wheel Slip}

The relationship between rotor speed and linear ground speed $v(t)$ accounts for transmission geometry and wheel slip:
\begin{equation}
    v(t) = \frac{\omega_m(t)}{i_g} r_w (1 - s(t))
\end{equation}
where $s(t) \in [0, 1)$ represents the slip ratio. Under nominal traction, $s(t) \approx 0$. When wheel slippage occurs, $s(t) \to 1$, which decouples wheel speed from chassis velocity ($v_{\mathrm{wheel}} \gg v_{\mathrm{ground}}$) while collapsing the ground reaction torque ($\tau_L \to 0$).

\subsection{Thermal Dissipation Models}

Winding resistance increases with copper temperature:
\begin{equation}
    R_a(T_m) = R_{a0} \left[ 1 + \alpha (T_m(t) - T_0) \right]
\end{equation}
where $R_{a0}$ is the baseline resistance at temperature $T_0$, and $\alpha$ is the copper temperature coefficient.

The thermal dynamics of the motor are modeled as a lumped first-order node accounting for Joule heating, core iron losses (Steinmetz formulation), and convective heat transfer to the ambient environment ($T_{\mathrm{amb}}$):
\begin{equation}
    C_{\mathrm{th}} \frac{dT_m(t)}{dt} = I_m(t)^2 R_a(T_m) + C_{fe} \omega_m(t)^{1.6} B_{\max}^2 - \frac{T_m(t) - T_{\mathrm{amb}}}{R_{\mathrm{th}}}
\end{equation}
Similarly, the battery pack temperature $T_{\mathrm{bat}}(t)$ accounts for internal resistive heating:
\begin{equation}
    C_{\mathrm{th,bat}} \frac{dT_{\mathrm{bat}}(t)}{dt} = I_{\mathrm{bat}}(t)^2 R_{\mathrm{int}} - \frac{T_{\mathrm{bat}}(t) - T_{\mathrm{amb}}}{R_{\mathrm{th,bat}}}
\end{equation}

\subsection{Battery Equivalent Circuit Model (ECM)}

Battery terminal voltage $V_{\mathrm{bat}}(t)$ is represented by an ohmic equivalent circuit:
\begin{equation}
    V_{\mathrm{bat}}(t) = V_{\mathrm{ocv}}(SOC) - I_{\mathrm{bat}}(t) R_{\mathrm{int}}(SOC, T_{\mathrm{bat}})
    \label{eq:battery_ecm}
\end{equation}
where open-circuit voltage $V_{\mathrm{ocv}}(SOC)$ scales with state of charge ($SOC$), and internal resistance $R_{\mathrm{int}}$ varies with $SOC$, temperature, and battery degradation. State of charge evolves via Coulomb counting based on total drawn power and charging currents:
\begin{equation}
    \frac{dSOC(t)}{dt} = -\frac{I_{\mathrm{bat}}(t)}{Q_{\mathrm{nom}}} \cdot 100
\end{equation}

\paragraph{Omission of the polarization branch ($V_{\mathrm{pol}}$):}
A complete 1-RC ECM introduces a capacitive polarization branch ($\frac{dV_{\mathrm{pol}}}{dt} = -\frac{V_{\mathrm{pol}}}{R_p C_p} + \frac{I_{\mathrm{bat}}}{C_p}$). This branch is deliberately omitted, and its transfer resistance $R_p$ is folded into $R_{\mathrm{int}}$ ($0.035 + 0.015 = 0.050\,\Omega$). 
\emph{Justification:} Detailed $R_p/C_p$ parameters cannot be calibrated without dedicated electrochemical pulse tests, and the polarization branch accounts for less than 20\% of total voltage sag. In a purely ohmic model, the load-step response is directly proportional to $I_{\mathrm{bat}} \cdot R_{\mathrm{int}}$, producing a clean, unambiguous signature for battery state-of-health (SOH) degradation without masking other anomalies.

\subsection{Physical Fault Injection and Transducer Emulation}

Rather than injecting synthetic noise into high-level features, faults perturb the physical equations directly:
\begin{itemize}
    \item \textbf{Mechanical Stall:} Imposes $\omega_m = 0$, collapsing back-EMF ($E_b = 0$) and causing the current to surge to the drive limit, accompanied by rapid thermal rise and battery voltage sag.
    \item \textbf{Wheel Slippage:} Forces $s(t) \to 1$, reducing load torque ($\tau_L \to 0$) while motor speed remains high.
    \item \textbf{Gearbox Degradation:} Injects an additive friction torque $\tau_{\mathrm{friction}}$, increasing steady-state current and operating temperature.
    \item \textbf{Stator Fault:} Injects harmonic current ripple and multiplies copper Joule heating.
    \item \textbf{Battery SOH Degradation:} Multiplies $R_{\mathrm{int}}$, producing pronounced voltage drops during acceleration steps.
\end{itemize}

Finally, physical states are conditioned through transducer models (Gaussian noise, ADC quantization, and thermal/gain drift) to generate realistic diagnostic telemetry for anomaly detection algorithms.

\subsection{Synthesized Telemetry Channels and Payload Mapping}
\label{sec:synthetic_dataset_generation}

The physical states computed by the model are serialized into the extended VDA~5050 telemetry payload published by each simulated vehicle. These sensor channels are grouped in a single \texttt{information} element of type \texttt{motor\_telemetry}, whose keys are emitted in a fixed order that forms part of the downstream data contract. The extended telemetry channels incorporated into the simulator are:

\begin{itemize}
    \item Motor current ($I_m$) (\texttt{motor\_current})
    \item Motor voltage ($V_m$) (\texttt{motor\_voltage})
    \item Motor temperature ($T_m$) (\texttt{motor\_temp})
    \item Motor angular velocity ($\omega_m$) (\texttt{motor\_speed\_rad\_s})
    \item Motor torque ($\tau_m$) (\texttt{motor\_torque\_nm})
    \item Back-electromotive force ($E_b$) (\texttt{back\_emf\_v})
    \item Wheel slip ratio ($s$) (\texttt{wheel\_slip})
    \item Motor power ($P_m$) (\texttt{motor\_power\_w})
    \item Battery current ($I_{bat}$) (\texttt{battery\_current\_a})
    \item Battery temperature ($T_{bat}$) (\texttt{battery\_temp\_c})
    \item Active fault flag (\texttt{fault\_active})
    \item Active fault types (\texttt{fault\_types})
\end{itemize}

These variables were selected because they are commonly used in predictive maintenance, condition monitoring, and fault diagnosis applications for electrically driven industrial vehicles. The two \texttt{fault\_*} channels carry the ground-truth identifiers of any injected anomaly, which is useful as a weak label for the detection pipeline. Battery state of charge, state of health and terminal voltage are intentionally not duplicated in this list: they continue to be reported through the standard VDA~5050 \texttt{batteryState} object. 

\subsection{AI-Assisted Implementation and Validation Workflow}
\label{subsec:ai_assisted_workflow}

The extension of the original \texttt{vda5050\_vehicle\_simulator} required the implementation of additional electromechanical, thermal, and battery modelling capabilities based on the physical framework described in the preceding subsections. Since the simulator adaptation itself was not the primary research objective of this work, modern AI-assisted software engineering techniques were employed to accelerate development. 

Specifically, an on-premise local AI model provided by the enterprise, NOVA, was primarily utilized during the design and iterative implementation phases, complemented by Google Gemini 3.8 Flash for architectural specification, code validation, and prompt structuring. The models were provided with detailed requirements regarding the desired physical behaviour, software architecture, VDA 5050 compatibility constraints, and telemetry channels to be incorporated into the simulator. 

The prompts were formulated as rigorous software specification documents detailing the required electromechanical, thermal, and battery equations, architectural constraints within the asynchronous Tokio runtime, Serde serialization structures, runtime fault scenarios triggered via VDA 5050 Instant Actions, and unit test requirements. 

Crucially, the AI-generated code was not incorporated directly into the project without verification. Instead, all algorithmic proposals were systematically reviewed, manually adapted to fit the simulator codebase, validated against theoretical physics curves, and thoroughly tested before final integration. A comprehensive validation of the simulator's physical responses and fault injection behavior is presented in Chapter~\ref{chap:testing_evaluation}.

\section{MQTT Broker Implementation and Integration} 
\label{sec:mqtt_broker_implementation}
\textit{Architectural counterpart: Chapter~\ref{chap:system_architecture}, Section~\ref{sec:arch_mqtt}.}

In the proposed architecture, Message Queuing Telemetry Transport (MQTT) serves as the lightweight edge messaging protocol connecting autonomous mobile robots to the data platform. As specified by the VDA~5050 standard, vehicle state and order messages are published over MQTT topics. To operationalize this layer within the containerized platform, the broker is deployed as a lightweight, containerised ingress service. The companion Telegraf bridge that forwards this edge traffic to the central event streaming backbone is described separately in Section~\ref{sec:stream_integration_implementation}.

\subsection{Eclipse Mosquitto Deployment and Broker Configuration}
\label{subsec:mosquitto_implementation}

The MQTT broker is implemented using Eclipse Mosquitto, deployed as a containerized service (\texttt{mosquitto}) on the dedicated Docker bridge network \texttt{mqtt-net}. The broker exposes standard port 1883 for unencrypted edge communication within the container environment.

Listing~\ref{lst:mosquitto-conf} illustrates the broker configuration (\texttt{mosquitto.conf}). Persistence is disabled (\texttt{persistence false}) because message durability and fault tolerance are delegated to the central event streaming log (Redpanda). All log output is routed to \texttt{stdout} to integrate with Docker's logging drivers.

\begin{lstlisting}[language=configToml, caption={Mosquitto configuration (mosquitto.conf)}, label={lst:mosquitto-conf}, basicstyle=\ttfamily\small, breaklines=true]
persistence false
log_dest stdout

listener 1883 0.0.0.0
allow_anonymous true
\end{lstlisting}

To prevent race conditions during stack initialization, a container health check is established in Listing~\ref{lst:mosquitto-docker-compose}. The probe uses \texttt{mosquitto\_pub} to publish a test message to a diagnostic topic every 5 seconds. Upstream telemetry producers (such as the vehicle simulator) and downstream consumers (such as Telegraf) utilize Docker's \texttt{condition: service\_healthy} dependency directive to ensure the broker is fully responsive before attempting connections.

\begin{lstlisting}[language=yaml, caption={Mosquitto Docker Compose configuration}, label={lst:mosquitto-docker-compose}, basicstyle=\ttfamily\small, breaklines=true]
services:
  mosquitto:
    image: eclipse-mosquitto:latest
    ports:
      - "1883:1883"
    volumes:
      - ./mosquitto/config:/mosquitto/config
    healthcheck:
      test: ["CMD", "mosquitto_pub", "-h", "localhost", "-p", "1883", "-t", "healthcheck", "-m", "1"]
      interval: 5s
      timeout: 5s
      retries: 10
    networks:
      - mqtt-net
\end{lstlisting}

\section{Event Stream Layer Implementation}
\label{sec:event_stream_implementation}
\textit{Architectural counterpart: Chapter~\ref{chap:system_architecture}, Section~\ref{sec:arch_event_streaming}.}

The event streaming layer was implemented using Redpanda, a Kafka-compatible
event streaming platform that acts as the central event bus connecting the
telemetry processing components. Three services are defined in the Docker
Compose configuration: the Redpanda broker, Redpanda Console for stream
monitoring and visualization, and a one-shot initialization service for
creating the required topics.

The Redpanda broker is configured to expose two Kafka-compatible endpoints.
The internal endpoint (\texttt{redpanda:9092}) enables communication between
services within the Docker network, while the external endpoint
(\texttt{localhost:19092}) allows access from the host machine. The
administrative API is exposed on port \texttt{9644}, and persistent storage
is configured through a Docker volume to preserve broker data across
container restarts.

To improve startup reliability, a health check is configured using
\texttt{rpk cluster health}. The Redpanda Console service waits until the
broker is healthy before starting, while the initialization service uses
the same dependency condition to ensure that topic creation is attempted
only after the broker is operational. The initialization service creates
the \texttt{vda5050-state}, \texttt{vda5050-state-deduplicated}, \texttt{fleet-features}, and \texttt{ai-anomalies} topics with four partitions and exits after
completing the operation. The \texttt{--if-not-exists} option prevents
errors when the topic has already been created.

A four-partition configuration was selected to provide parallelism for
the expected telemetry workload. Depending on message size, hardware
resources, replication settings, and processing overhead, a partition
may support throughput in the order of thousands or tens of thousands
of messages per second. Therefore, four partitions provide a reasonable
initial configuration for the simulated fleet, although the effective
capacity must be validated through performance testing. The broker runs
in \texttt{dev-container} mode with one CPU core assigned through
\texttt{--smp 1}, reflecting the resource constraints of the local
development environment.

\begin{lstlisting}[language=yaml, caption={Redpanda Docker Compose}, label={lst:redpanda-docker-compose}, basicstyle=\ttfamily\small, breaklines=true]
services:
  redpanda:
    image: docker.redpanda.com/redpandadata/redpanda:v26.1.12
    hostname: redpanda
    command:
      - redpanda
      - start
      - --kafka-addr internal://0.0.0.0:9092,external://0.0.0.0:19092
      - --advertise-kafka-addr internal://redpanda:9092,external://localhost:19092
      - --rpc-addr redpanda:33145
      - --advertise-rpc-addr redpanda:33145
      - --mode dev-container
      - --smp 1
      - --default-log-level=info
    ports:
      - 19092:19092
      - 19644:9644
    volumes:
      - redpanda-data:/var/lib/redpanda/data
    networks:
      - redpanda-net
    healthcheck:
      test: ["CMD-SHELL", "rpk cluster health --exit-when-healthy"]
      interval: 10s
      timeout: 5s
      retries: 10
      start_period: 30s

  redpanda-console:
    image: docker.redpanda.com/redpandadata/console:v3.8.0
    entrypoint: /bin/sh
    command: -c 'echo "$$CONSOLE_CONFIG_FILE" > /tmp/config.yml; /app/console'
    environment:
      CONFIG_FILEPATH: /tmp/config.yml
      CONSOLE_CONFIG_FILE: |
        kafka:
          brokers: ["redpanda:9092"]
        redpanda:
          adminApi:
            enabled: true
            urls: ["http://redpanda:9644"]
    ports:
      - 9080:8080
    depends_on:
      redpanda:
        condition: service_healthy
    networks:
      - redpanda-net

  redpanda-init:
    image: docker.redpanda.com/redpandadata/redpanda:v26.1.12
    restart: "no"
    entrypoint: ["/bin/sh", "-c"]
    command:
      - |
        rpk topic create vda5050-state --partitions 4 -X brokers=redpanda:9092 --if-not-exists
        rpk topic create fleet-features --partitions 4 -X brokers=redpanda:9092 --if-not-exists
        rpk topic create ai-anomalies --partitions 4 -X brokers=redpanda:9092 --if-not-exists
    depends_on:
      redpanda:
        condition: service_healthy
    networks:
      - redpanda-net

volumes:
  redpanda-data:

\end{lstlisting}

\section{Streaming Integration and Inter-Service Connectivity}
\label{sec:stream_integration_implementation}
\textit{Architectural counterpart: Chapter~\ref{chap:system_architecture}, Section~\ref{sec:arch_streaming_integration}.}

This section consolidates the connector services that bridge the different architectural tiers, mirroring the integration topology described in Chapter~\ref{chap:system_architecture}. Two bridges are implemented: the Telegraf agent that forwards edge MQTT telemetry into Redpanda, and the Kafka Connect cluster that sinks enriched streams from Redpanda into QuestDB.

\subsection{Telemetry Bridging via Telegraf}
\label{subsec:telegraf_integration}
\textit{Architectural counterpart: Chapter~\ref{chap:system_architecture}, Section~\ref{sec:arch_streaming_integration}.}

A key architectural requirement is the clean decoupling between the edge MQTT protocol and the central Redpanda streaming tier. Because the MQTT boundary only requires protocol adaptation between MQTT and Kafka-compatible streams, Telegraf was selected as a dedicated telemetry ingestion agent.

Telegraf operates as a dual-homed container attached simultaneously to \texttt{mqtt-net} and \texttt{redpanda-net}. This network topology enforces strict isolation: source systems and the MQTT broker remain in \texttt{mqtt-net} with zero direct exposure to the internal Kafka network, while Telegraf transparently forwards incoming payloads across the boundary.

Listing~\ref{lst:telegraf-docker-compose} shows the Docker Compose configuration for the Telegraf bridging service.

\begin{lstlisting}[language=yaml, caption={Telegraf Docker Compose configuration}, label={lst:telegraf-docker-compose}, basicstyle=\ttfamily\small, breaklines=true]
services:
  telegraf:
    image: telegraf:1.39-alpine
    restart: unless-stopped
    command: telegraf --config /etc/telegraf/conf/telegraf.conf
    volumes:
      - ./telegraf.conf:/etc/telegraf/conf/telegraf.conf:ro
    environment:
      MQTT_BROKER_URL: ${MQTT_BROKER_URL:-tcp://mosquitto:1883}
      KAFKA_BROKER: ${TELEGRAF_KAFKA_BROKER:-redpanda:9092}
    networks:
      - mqtt-net
      - redpanda-net
\end{lstlisting}

The bridging logic is governed by three coordinated stages configured in \texttt{telegraf.conf} (Listing~\ref{lst:telegraf-conf}):
\begin{enumerate}
    \item \textbf{MQTT Telemetry Ingestion (\texttt{inputs.mqtt\_consumer}):} Telegraf subscribes to the VDA~5050 state topic hierarchy using the wildcard pattern \texttt{uagv/v2/+/+/state}, matching all vehicles across arbitrary manufacturers. Ingestion uses Quality of Service 1 (QoS~1) with persistent sessions, preventing telemetry loss during momentary network disconnects. Telemetry payloads are ingested as raw string values (\texttt{data\_format = "value"}) to prevent premature deserialization.
    \item \textbf{Topic Parsing and Composite Key Generation (\texttt{processors.template}):} The hierarchical MQTT topic structure is parsed into semantic tags:
    \begin{quote}
        \texttt{\_/\_/source\_manufacturer/source\_serial\_number/message\_type}
    \end{quote}
    A processor template constructs a composite key \texttt{source\_asset\_key} formatted as \texttt{<manufacturer>:<serial\_number>} (e.g., \texttt{rikeb:s11}). This composite key serves as the Kafka message routing key (\texttt{routing\_tag}), guaranteeing two crucial properties in Redpanda:
    \begin{itemize}
        \item \emph{Collision avoidance:} AGVs from different manufacturers with identical serial numbers are unambiguously partitioned.
        \item \emph{Chronological partition ordering:} Messages from the same physical asset consistently hash to the same Kafka partition, ensuring strict temporal order for downstream state reconstruction.
    \end{itemize}
    \item \textbf{Unmutated Kafka Publishing and Header Enrichment (\texttt{outputs.kafka}):} Telegraf publishes to the centralized \texttt{vda5050-state} topic in Redpanda. Crucially, the VDA~5050 JSON payload is emitted completely unmodified (\texttt{template = '\{\{ .Field "value" \}\}'}). Schema transformations and data validation are intentionally deferred to downstream Spark Structured Streaming jobs, preserving the raw data fidelity required for the Bronze lakehouse tier. Simultaneously, Telegraf appends enriched Kafka record headers (\texttt{source\_serial\_number}, \texttt{source\_manufacturer}, \texttt{source\_asset\_key}, \texttt{message\_type}, and \texttt{mqtt\_topic}), providing full protocol provenance to consumers without altering message bodies.
\end{enumerate}

\begin{lstlisting}[language=configToml, caption={Telegraf bridging configuration (telegraf.conf)}, label={lst:telegraf-conf}, basicstyle=\ttfamily\small, breaklines=true]
[agent]
  interval = "1s"
  flush_interval = "1s"
  precision = "1ms"
  omit_hostname = true

[[inputs.mqtt_consumer]]
  servers = ["${MQTT_BROKER_URL}"]
  topics = ["uagv/v2/+/+/state"]
  client_id = "telegraf-vda5050"
  qos = 1
  persistent_session = true
  data_format = "value"
  data_type = "string"

  [[inputs.mqtt_consumer.topic_parsing]]
    topic = "uagv/v2/+/+/state"
    tags = "_/_/source_manufacturer/source_serial_number/message_type"

[[processors.template]]
  order = 0
  tag = "source_asset_key"
  template = '{{ .Tag "source_manufacturer" }}:{{ .Tag "source_serial_number" }}'

[[outputs.kafka]]
  brokers = ["${KAFKA_BROKER}"]
  topic = "vda5050-state"
  routing_tag = "source_asset_key"
  required_acks = -1
  idempotent_writes = true
  compression_codec = 4
  startup_error_behavior = "retry"
  data_format = "template"
  template = '{{ .Field "value" }}'

  [outputs.kafka.headers]
    source_serial_number = '{{ index .Tags "source_serial_number" }}'
    source_manufacturer = '{{ index .Tags "source_manufacturer" }}'
    source_asset_key = '{{ index .Tags "source_asset_key" }}'
    message_type = '{{ index .Tags "message_type" }}'
    mqtt_topic = '{{ index .Tags "topic" }}'
\end{lstlisting}

\subsection{Kafka Connect Integration and Custom Single Message Transformations}
\label{subsec:questdb-kafka-connect}
\textit{Architectural counterpart: Chapter~\ref{chap:system_architecture}, Section~\ref{sec:arch_streaming_integration}.}

To transfer telemetry seamlessly from the Redpanda broker into QuestDB, a specialized Kafka Connect service (\texttt{questdb-kafka-connect}) was implemented. Because QuestDB is an append-only relational time-series database that expects tabular scalar structures, directly ingesting complex, hierarchically nested VDA~5050 JSON payloads (containing nested objects and arrays) poses technical challenges.

\textit{(TO DO: SINK CONNECTOR IMPLEMENTATION)}

\section{Stream Processing Layer: Apache Spark Implementation}
\label{sec:spark-implementation}
\textit{Architectural counterpart: Chapter~\ref{chap:system_architecture}, Section~\ref{sec:spark-stream-processing}.}
The stream processing tier is implemented using \textbf{Apache Spark} (version 3.5.0) via its \textbf{Structured Streaming} framework. It is deployed as a dedicated containerized service (\texttt{spark-processing}) within the Docker Compose orchestration stack.


\subsection{Docker Compose Service Definition}
\label{subsec:spark-docker-definition}

Listing~\ref{lst:spark-docker-compose} presents the Docker Compose definition of the streaming tier. 

\paragraph{Initialization:}
The container initialization is governed by health-check and completion dependencies:
\begin{itemize}
    \item It waits for the \texttt{redpanda} broker to report a healthy cluster status.
    \item It awaits the successful execution of \texttt{redpanda-init} to ensure that all required topics (including \texttt{vda5050-state} and \texttt{vda5050-state-deduplicated}) have been provisioned before job submission.
    \item It requires \texttt{minio} to be started and the \texttt{createbuckets} utility container to successfully complete bucket initialization (\texttt{lakehouse}).
\end{itemize}

Once all dependencies are satisfied, streaming execution is launched using \texttt{spark-submit} in local client mode (\texttt{local[*]}), dynamically resolving the required Kafka integration library (\texttt{spark-sql-kafka-0-10}) and Amazon S3A connector (\texttt{hadoop-aws}).

\begin{lstlisting}[language=yaml, caption={Apache Spark Docker Compose configuration}, label={lst:spark-docker-compose}, basicstyle=\ttfamily\small, breaklines=true]
services:
  spark-processing:
    image: apache/spark:3.5.0
    container_name: spark-processing
    depends_on:
      redpanda:
        condition: service_healthy
      redpanda-init:
        condition: service_completed_successfully
      minio:
        condition: service_started
      createbuckets:
        condition: service_completed_successfully
    networks:
      - redpanda-net
      - lakehouse-net
    volumes:
      - ./jobs:/opt/spark/jobs
    environment:
      - KAFKA_BOOTSTRAP_SERVERS=${KAFKA_BROKER:-redpanda:9092}
      - KAFKA_INPUT_TOPIC=${KAFKA_INPUT_TOPIC:-vda5050-state}
      - S3_ENDPOINT=${S3_ENDPOINT:-http://minio:9000}
      - S3_ACCESS_KEY=${S3_ACCESS_KEY:-admin}
      - S3_SECRET_KEY=${S3_SECRET_KEY:-admin123456}
      - S3_CHECKPOINT_BASE=${S3_CHECKPOINT_BASE:-s3a://lakehouse/checkpoints}
      - ICEBERG_CATALOG=${ICEBERG_CATALOG:-lakehouse}
      - ICEBERG_WAREHOUSE=${ICEBERG_WAREHOUSE:-s3a://lakehouse/iceberg}
      - ICEBERG_BRONZE_TABLE=${ICEBERG_BRONZE_TABLE:-lakehouse.bronze.vda5050_state}
      - PROCESSING_INTERVAL=${PROCESSING_INTERVAL:-15 seconds}
    command: >
      /opt/spark/bin/spark-submit
      --conf spark.jars.ivy=/tmp/.ivy2
      --packages org.apache.spark:spark-sql-kafka-0-10_2.12:3.5.0,org.apache.hadoop:hadoop-aws:3.3.4,org.apache.iceberg:iceberg-spark-runtime-3.5_2.12:1.5.2
      --master local[*]
      /opt/spark/jobs/vda5050_state_writer.py

  spark-maintenance:
    image: apache/spark:3.5.0
    profiles:
      - maintenance
    depends_on:
      spark-processing:
        condition: service_started
    networks:
      - lakehouse-net
    volumes:
      - ./jobs:/opt/spark/jobs
    environment:
      - S3_ENDPOINT=${S3_ENDPOINT:-http://minio:9000}
      - S3_ACCESS_KEY=${S3_ACCESS_KEY:-admin}
      - S3_SECRET_KEY=${S3_SECRET_KEY:-admin123456}
      - ICEBERG_CATALOG=${ICEBERG_CATALOG:-lakehouse}
      - ICEBERG_WAREHOUSE=${ICEBERG_WAREHOUSE:-s3a://lakehouse/iceberg}
      - ICEBERG_BRONZE_TABLE=${ICEBERG_BRONZE_TABLE:-lakehouse.bronze.vda5050_state}
    command: >
      /opt/spark/bin/spark-submit
      --conf spark.jars.ivy=/tmp/.ivy2
      --packages org.apache.hadoop:hadoop-aws:3.3.4,org.apache.iceberg:iceberg-spark-runtime-3.5_2.12:1.5.2
      --master local[*]
      /opt/spark/jobs/maintenance_bronze.py

networks:
  redpanda-net:
  lakehouse-net:
\end{lstlisting}

\paragraph{Apache Iceberg Catalog Configuration:}
Metadata management in the streaming tier relies on a \textbf{Hadoop-based Iceberg catalog} rather than on an external catalog service.
It is instantiated through the Spark \texttt{SparkCatalog} implementation by setting \texttt{spark.sql.catalog.\allowbreak lakehouse.type = hadoop} and pointing its \texttt{warehouse} to the object-store location \texttt{s3a://lakehouse/iceberg}.
In this configuration the Iceberg metadata tree (table snapshots, schemas and manifests) is persisted directly as files inside the MinIO bucket through the Hadoop filesystem abstraction, which removes the need for a standalone metastore such as a Thrift or REST catalog and further reduces the number of operational dependencies.
Table discovery and atomic metadata commits are therefore resolved against the S3-compatible object store itself, in line with the Lakehouse design presented in Chapter~\ref{chap:system_architecture} (Section~\ref{subsec:table-format-iceberg}).

\subsection{Jobs} 
The streaming layer is implemented through a set of dedicated Spark Structured Streaming jobs, each responsible for a specific stage of the data pipeline. 
Three jobs compose the current architecture: 
  \begin{itemize} 
    \item \texttt{vda5050\_state\_writer}: Persists raw VDA5050 telemetry from Kafka into the Bronze layer. 
    \item \texttt{fleet\_features\_writer}: Extracts, transforms and enriches telemetry into a machine-learning-ready feature representation. 
    \item \texttt{ai\_anomalies\_writer}: Executes real-time anomaly detection and publishes anomaly scores for downstream consumption. 
  \end{itemize} The latter two jobs implement a Dual-Sink Output Pattern, simultaneously persisting processed data into the data lake and publishing the resulting stream back to Redpanda for further processing. 
  
  \paragraph{Dual-Sink Output Pattern} The feature engineering and anomaly detection jobs write their outputs to two independent destinations: 
    \begin{itemize} 
      \item \textbf{Analytical Persistence Sink (MinIO):} Processed datasets are stored as Apache Iceberg tables within the data lakehouse, enabling historical analytics, schema evolution and time-travel queries. 
      \item \textbf{Event Streaming Sink (Redpanda):} Processed records are republished to Kafka topics, allowing downstream consumers to access enriched telemetry and anomaly detection results in real time. \end{itemize}

\subsection{Bronze Ingestion Job} \label{subsec:bronze-ingestion-job} 
The Bronze ingestion job (\texttt{vda5050\_state\_writer.py}) constitutes the entry point of the Medallion architecture. 
Its responsibility is to persist the incoming telemetry stream without modification, creating an immutable historical record of the original VDA5050 messages.
Unlike subsequent processing stages, no feature extraction, filtering, aggregation or deduplication is applied during Bronze ingestion. 
The objective of this layer is to preserve the raw telemetry exactly as it was received from the messaging infrastructure, ensuring full data lineage, reproducibility and auditability. 
\paragraph{Ingestion and Storage Configuration} The job is implemented using PySpark Structured Streaming and continuously consumes records from the Kafka topic \texttt{vda5050-state}. 
Incoming messages are processed in micro-batches and persisted into an Apache Iceberg table stored in the MinIO object store.
Spark is configured to access MinIO through the S3A filesystem connector, using path-style access to the endpoint:\texttt{http://minio:9000}.

The target Iceberg table represents the Bronze layer of the data platform and stores the complete telemetry payload together with ingestion metadata required for historical analysis and time-travel operations.
\paragraph{Ingestion and Storage Configuration} 
The job executes micro-batches at a configurable interval, defined by the environment variable \texttt{PROCESSING\_INTERVAL}.
\subsection{Silver Feature Engineering Job}
Para la implementación de la capa de silver, se crea un job llamado, fleet-features-writer.py, el cual se encarga de flattening y de features extractions.
Lee del topic, vda5050-state, hace flattening, y feature extraction, y luego escribe en el topic fleet-features, y en la tabla de iceberg con el mismo nombre.
El job reutiliza las funciones de sink.py y de spark.
Para el tiempo de procesamiento, se utiliza un microbatch personalizable, el cual se puede configurar a través de la variable de entorno. 
En el caso de silver, se establece un tiempo de procesamiento por defecto menor al de bronze, debido al requisito de real time para fleet-features, que tiene que ser posteriormente en la parte de inferencia analizado en tiempo real.
\label{subsec:silver-feature-engineering-job}

\textit{(TO DO)}

\subsection{Gold Inference Job}
\label{subsec:gold-inference-job}

\textit{(TO DO)}

\section{Data Lakehouse Layer Implementation}
\label{sec:data-lakehouse-implementation}
\textit{Architectural counterpart: Chapter~\ref{chap:system_architecture}, Section~\ref{sec:data-lakehouse-layer}.}

\textit{(TO BE COMPLETED: Overview of the Data Lakehouse implementation, storage topology, and integration between object storage and modern table formats).}

\subsection{S3-Compatible Object Storage: MinIO}
\label{subsec:minio-s3-implementation}

The object storage tier was implemented using \textbf{MinIO}, an open-source, high-performance distributed object storage server fully compatible with the Amazon S3 API. In this architecture, MinIO serves as the primary persistence layer for storing raw Bronze Parquet files and managing streaming checkpoints for Apache Spark.

Two services are defined within the Docker Compose deployment:
\begin{itemize}
    \item \textbf{MinIO Server (\texttt{minio}):} Runs the core object storage engine, persisting data in a dedicated Docker volume (\texttt{minio-data}). The S3 API operates on container port \texttt{9000}, which is mapped to host port \texttt{19000} to prevent host port collisions with QuestDB. The interactive web console is exposed on host port \texttt{9001}. A container health check probes \texttt{/minio/health/live} to guarantee service availability before downstream consumers initiate requests.
    \item \textbf{Bucket Provisioning (\texttt{createbuckets}):} A lightweight initialization container using the MinIO Client (\texttt{minio/mc}) CLI. It waits for the MinIO server to report a healthy status, configures a local connection alias, and automatically provisions the target bucket (\texttt{lakehouse}) using \texttt{mc mb --ignore-existing}. Once bucket creation succeeds, the container terminates cleanly.
\end{itemize}

Both services are attached to the \texttt{lakehouse-net} network, ensuring secure, isolated network communication with the Spark stream processing tier.

\begin{lstlisting}[language=yaml, caption={MinIO Docker Compose configuration}, label={lst:minio-docker-compose}, basicstyle=\ttfamily\small, breaklines=true]
services:
  minio:
    image: minio/minio
    command: server /data --console-address ":9001"
    environment:
      MINIO_ROOT_USER: admin
      MINIO_ROOT_PASSWORD: admin123456
    volumes:
      - minio-data:/data
    ports:
      - "19000:9000"
      - "9001:9001"
    networks:
      - lakehouse-net
    healthcheck:
      test: ["CMD", "curl", "-f", "http://localhost:9000/minio/health/live"]
      interval: 10s
      timeout: 5s
      retries: 5

  createbuckets:
    image: minio/mc
    depends_on:
      minio:
        condition: service_healthy
    entrypoint: >
      /bin/sh -c "
      sleep 5;
      mc alias set local http://minio:9000 admin admin123456;
      mc mb --ignore-existing local/lakehouse;
      exit 0;
      "
    networks:
      - lakehouse-net

volumes:
  minio-data:
\end{lstlisting}

\subsection{Table Format Implementation: Apache Iceberg}
\label{subsec:iceberg-implementation}

\textit{(TO BE COMPLETED: Apache Iceberg catalog configuration, table schema definition, and partitioning strategy across Bronze, Silver, and Gold tiers).}

\section{QuestDB Implementation and Integration}
\label{sec:questdb-implementation}
\textit{Architectural counterpart: Chapter~\ref{chap:system_architecture}, Section~\ref{subsec:bifurcated-storage}.}

As established in the architectural design (Chapter~\ref{chap:system_architecture}), the platform adopts a bifurcated persistence strategy that strictly separates the analytical lakehouse storage (MinIO and Apache Iceberg) from the operational time-series database. While the lakehouse is optimized for heavy, multi-tier batch analytics and offline model training, real-time observability requires a database capable of sub-second streaming ingestion and high-frequency SQL aggregations across sliding time windows. To fulfill this operational role, \textbf{QuestDB} (version 9.4.3) was selected and integrated into the containerized platform.

The QuestDB subsystem interfaces directly with the central event streaming layer through a dedicated, containerized \textbf{Apache Kafka Connect} cluster. This design establishes an automated ingestion pipeline that streams live VDA~5050 robot telemetry from Redpanda into QuestDB without requiring intermediate message-polling services or bespoke application logic.

\subsection{Database Deployment and Schema Architecture}
\label{subsec:questdb-deployment-schema}

The database engine is deployed as an isolated service (\texttt{questdb}) within a dedicated bridge network (\texttt{questdb-net}). Persistent telemetry is stored on a dedicated Docker volume (\texttt{questdb-data}), preserving all historical state messages across container restarts.

QuestDB provides native multi-protocol support. In this deployment, three distinct interfaces are configured:
\begin{itemize}
    \item \textbf{HTTP REST API and Web Console (port 9000):} Used for system administration, automated schema provisioning, and diagnostic telemetry inspection.
    \item \textbf{PostgreSQL Wire Protocol (port 8812):} QuestDB natively emulates the PostgreSQL wire protocol. To adhere to the principle of least privilege, a dedicated read-only user (\texttt{QDB\_PG\_READONLY\_USER: grafana}) is configured. This endpoint allows Grafana dashboards to query operational tables directly using standard SQL without exposing write permissions.
    \item \textbf{InfluxDB Line Protocol (port 9009):} Maintained for low-overhead, high-throughput streaming metrics ingestion if required.
\end{itemize}

A container health check continuously probes the HTTP status endpoint (\texttt{http://localhost:9003/status}) to ensure the engine is fully initialized before downstream connectors attempt data transmission.

To automate database table provisioning upon initial startup, a one-shot container (\texttt{questdb-init}) is defined using \texttt{curlimages/curl}. It executes after QuestDB reports a healthy state, submitting SQL DDL commands to QuestDB's \texttt{/exec} REST endpoint.

\textit{(TO DO: TABLE IMPLEMENTATIONS)}

\subsection{Automated Lifecycle and Dependency Management}
\label{subsec:questdb-lifecycle}

To guarantee reliable system initialization, service orchestration is strictly ordered through Docker health checks and one-shot provisioning containers:
\begin{enumerate}
    \item The \texttt{questdb} service initializes its engine and validates readiness on port 9003.
    \item The \texttt{questdb-init} service triggers the SQL DDL creation of the \texttt{vda5050\_state} table.
    \item The \texttt{questdb-kafka-connect} service initializes only after the Redpanda broker is healthy, all required Kafka topics are created (\texttt{redpanda-init}), and QuestDB tables exist (\texttt{questdb-init}).
    \item Finally, a dedicated initialization container (\texttt{questdb-kafka-connect-init}) executes an automated shell script that submits the connector configuration via the Kafka Connect REST API (\texttt{PUT /connectors/<name>/config}). The script actively polls the connector status endpoint until the task achieves the \texttt{RUNNING} state, ensuring that the pipeline is completely operational before simulated AGVs begin streaming messages.
\end{enumerate}

This architecture guarantees end-to-end automated orchestration, preventing race conditions during platform startup while maintaining strict network isolation between operational storage, event streaming, and analytical tiers.

\textit{(TO DO: DOCKER COMPOSE)}

\section{Machine Learning Pipeline Implementation}
\label{sec:ml_pipeline_implementation}
\textit{Architectural counterpart: Chapter~\ref{chap:system_architecture}, Section~\ref{sec:arch_ml}.}

\textit{(TO DO)}

\section{Grafana Implementation}
\label{sec:grafana_implementation}
\textit{Architectural counterpart: Chapter~\ref{chap:system_architecture}, Section~\ref{sec:arch_observability}.}

\textit{(TO DO)}

\section{Security and Network Isolation Enforcement}
\label{sec:security_implementation}
\textit{Architectural counterpart: Chapter~\ref{chap:system_architecture}, Section~\ref{sec:arch_security}.}

\textit{(TO DO)}
